\uuid{eIUU}
\exo7id{1903}
\titre{exo7 1903}
\auteur{maillot}
\organisation{exo7}
\datecreate{2001-09-01}
\isIndication{false}
\isCorrection{false}
\chapitre{Suite}
\sousChapitre{Suite dans Rn}
\cpge{mpsi-mp2i.nombres-reels-suites.c, pcsi.nombres-reels-suites.b, ptsi.nombres-reels-suites.b, mpsi-mp2i.nombres-reels-suites.d, pcsi.nombres-reels-suites.c, ptsi.nombres-reels-suites.c, mpsi-mp2i.nombres-reels-suites.f, pcsi.nombres-reels-suites.e, ptsi.nombres-reels-suites.e, mpsi-mp2i.limites-continuite.b, pcsi.limites-continuite.b, ptsi.limites-continuite.b, bcpst1.suites-reelles, bcpst1.limites-continuite.c, mp-mpi.topologie-espaces-normes.b, pc.espaces-vectoriels-normes.b, psi.espaces-vectoriels-normes.b, mp-mpi.topologie-espaces-normes.e, pc.espaces-vectoriels-normes.e, psi.espaces-vectoriels-normes.e, pt.fonctions-deux-trois-variables.b}
\module{Analyse}
\niveau{L1}
\difficulte{}

\contenu{
\question{
Soit $f:\Rr^d\rightarrow\Rr^d$ continue.
Soit $x_0\in\Rr^d$. Soit $x_n$ la suite d\'efinie par
$$x_{n+1}=f(x_n).$$
Supposons que $||x_n-x_{n+1}||\rightarrow 0$.
Montrer que si $a\in\ A$ alors $f(a)=a$.

Indication: appliquer la d\'efinition de la continuit\'e
de $f$ en $a$ en termes de limites.
}
}
